%\subsection{Versuchsdurchf�hrung}
%Ein Praktikumsbericht sowie andere wissenschaftliche Arbeiten unterliegen einer Gliederung aus konkreten Elementen.
%Die Ausarbeitung, die w�hrend den Praxisphasen an der HSHL angefertigt werden muss, wird ebenfalls als Praktikumsbericht angesehen.
%Laut Balzert (\cite{Balz2008}, S. 154) enth�lt ein Praktikumsbericht folgende Bestandteile:
%
%\begin{itemize}
%\item Deckblatt
%\item Sperrvermerk (evtl.)
%\item Inhaltsverzeichnis
%\item Textteil bestehend aus Einleitung, Hauptteil und Schluss
%\item Anhang (falls erforderlich)
%\item Literaturverzeichnis (evtl.)
%\item Eidesstattliche Versicherung (falls erforderlich)
%\end{itemize}
%
%\noindent
%Die Eidesstaatliche Versicherung, die auch diesem Dokument beigef�gt wurde, ist auf eine gesonderte Seite am Ende der Bachelorarbeit zu setzen.
%Eine Projektarbeit, die w�hrend des Praxissemesters im Unternehmen geschrieben wird, enth�lt neben den bereits genannten Elementen ein Literaturverzeichnis und einen Sperrvermerk.
%Sobald mehr als f�nf Abbildungen, Abk�rzungen oder Symbole in der Ausarbeitung verwendet werden, sollte ein entsprechendes Verzeichnis angelegt werden.
%Das Abk�rzungs-, Abbildungs-, Tabellen-, Formel und Symbolverzeichnis steht am Anfang der der wissenschaftlichen Ausarbeitung direkt hinter dem Inhaltsverzeichnis.
%F�r die Bachelor- sowie Masterarbeit werden nach Balzert (\cite{Balz2008}, S. 154) folgende zus�tzliche Bestandteile empfohlen:\footnote{In dieser wissenschaftlichen Arbeit sind alle erforderlichen formalen Bestandteile einer Bachelorarbeit enthalten.}
%
%\begin{itemize}
%\item Abk�rzungsverzeichnis
%\item Abbildungs- und Tabellenverzeichnis
%\item Formel- und Symbolverzeichnis (falls erforderlich)
%\end{itemize}
%
%\noindent
%�blicherweise umfasst der Textumfang eines Praktikumsberichts 5 bis 15 Textseiten.
%Hausarbeiten bewegen sich in einem Rahmen von 15 bis 30.
%Bei Projektarbeiten wird je nach Thematik und Umfang ein Textumfang von 10 bis 50 Seiten gefordert.
%Eine Bachelor- oder Masterarbeit sollte einen reinen Textanteil von mindestens 30 bis 80 Seiten besitzen.
%Als Textumfang wird die Anzahl der Seiten von Einleitung, Hauptteil und Schluss und nicht die Anzahl der Seiten insgesamt angesehen (vgl. \cite{Balz2008}, S. 152).
